\section{Anexo D Docker}\label{anexo-d-docker}

\subsection{ANEXO D.1. CONTENEDORIZACIÓN E INFRAESTRUCTURA DE DESPLIEGUE
(DOCKER)}\label{anexo-d.1.-contenedorizaciuxf3n-e-infraestructura-de-despliegue-docker}

En el ámbito del desarrollo de software moderno y la ciberseguridad, la
reproducibilidad del entorno de ejecución constituye un pilar crítico.
Tradicionalmente, el despliegue de aplicaciones que integran múltiples
herramientas de análisis (como compiladores de Solidity y motores de
ejecución simbólica) se enfrentaba al problema de ``funciona en mi
máquina'', derivado de las discrepancias en las versiones de las
dependencias, librerías del sistema operativo y configuraciones locales.
Para mitigar este riesgo, el presente proyecto adopta una arquitectura
basada en contenedores de aplicación a través del ecosistema de
\textbf{Docker} y \textbf{Docker Compose}.

\subsection{ANEXO D.2. Fundamentos de Contenedorización: Docker y Docker
Compose}\label{anexo-d.2.-fundamentos-de-contenedorizaciuxf3n-docker-y-docker-compose}

Docker es una plataforma de código abierto basada en la tecnología de
contenedorización, la cual permite empaquetar una aplicación y todas sus
dependencias (binarios, librerías, archivos de configuración) en una
unidad estandarizada denominada \textbf{contenedor}. A diferencia de la
virtualización tradicional, Docker opera mediante la virtualización a
nivel de sistema operativo, compartiendo el núcleo (kernel) del sistema
anfitrión pero ejecutando los procesos en espacios de usuario
completamente aislados a través de namespaces y cgroups. Desde la
perspectiva de la seguridad, este aislamiento garantiza que los procesos
del pipeline de auditoría de EVMAudit se ejecuten de forma confinada,
mitigando el impacto en la infraestructura anfitriona ante la eventual
ejecución de código arbitrario o inesperado durante el análisis de
contratos inteligentes.

Por su parte, \textbf{Docker Compose} es la herramienta diseñada para
definir y orquestar aplicaciones Docker multi-contenedor. Mediante el
uso de un archivo de configuración declarativo en formato YAML
(docker-compose.yml), permite definir con precisión los servicios que
componen el sistema, sus dependencias de arranque, la exposición de
puertos hacia el exterior, la creación de redes aisladas y la asignación
de volúmenes persistentes. En el contexto de EVMAudit, actúa como el
motor de despliegue unificado, permitiendo al administrador inicializar
toda la infraestructura del TFM (servidor FastAPI, interfaz web y
almacenamiento de informes) de manera centralizada.

\subsection{ANEXO D.3. Estrategia de Construcción de la Imagen
(Dockerfile)}\label{anexo-d.3.-estrategia-de-construcciuxf3n-de-la-imagen-dockerfile}

La construcción de la imagen se define en un único Dockerfile
optimizado. Debido a que el pipeline de análisis requiere interactuar
con el sistema operativo para invocar compiladores y binarios de
seguridad, se ha seleccionado \textbf{Ubuntu 22.04} como imagen base,
proporcionando un entorno estable y con soporte extendido para
dependencias nativas de Linux en arquitectura amd64.

El proceso de aprovisionamiento de la imagen se divide en las siguientes
fases críticas: 1. \textbf{Entorno y Variables de Sistema}: Se
configuran las variables de entorno \texttt{PYTHONDONTWRITEBYTECODE=1} y
\texttt{PYTHONUNBUFFERED=1} para optimizar la ejecución de Python dentro
del contenedor, evitando la escritura de residuos binarios y forzando el
volcado de logs en tiempo real. Asimismo, se establece
\texttt{DEBIAN\_FRONTEND=noninteractive} para suprimir diálogos
interactivos durante la instalación de paquetes. 2.
\textbf{Aprovisionamiento de Compiladores (Solidity):} Se añade el
repositorio PPA oficial de Ethereum (ppa:ethereum/ethereum) para
incorporar el compilador nativo de Solidity (solc). Posteriormente, se
instala la utilidad \texttt{solc-select} mediante el gestor de paquetes
de Python para automatizar la descarga y conmutación de versiones. 3.
\textbf{Integración del Fuzzer Echidna:} Dado que Echidna se distribuye
de manera óptima como un binario estático para Linux, el contenedor
automatiza su descarga directa (versión v2.3.2) desde los repositorios
oficiales de \emph{Crytic}, procediendo a su extracción e instalación en
\texttt{/usr/local/bin/} para garantizar su disponibilidad inmediata en
el PATH del sistema. 4. \textbf{Optimización de Dependencias con
\texttt{uv} (Multi-stage Build):} Con el objetivo de minimizar los
tiempos de construcción y asegurar una gestión eficiente de los paquetes
de Python, se emplea un mecanismo de construcción en etapas múltiples
(Multi-stage build), importando los binarios optimizados del gestor uv
directamente desde su imagen oficial en el registro de GitHub
(ghcr.io/astral-sh/uv:latest). 5. \textbf{Instalación del Paquete
Local:} Tras establecer el directorio de trabajo en /app y copiar el
código fuente , se ejecuta el comando
\texttt{uv\ sync\ -\/-frozen\ -\/-no-cache}. Esto resuelve de forma
determinista el grafo de dependencias del archivo uv.lock, registrando
el paquete local editable evmaudit e instalando el servidor ASGI Uvicorn
sin almacenar datos residuales en la caché de la imagen.

\subsection{ANEXO D.4. Orquestación de Servicios (Docker
Compose)}\label{anexo-d.4.-orquestaciuxf3n-de-servicios-docker-compose}

La coordinación del contenedor web y sus dependencias con el sistema
anfitrión se gestiona de forma declarativa mediante un archivo
docker-compose.yml. La especificación del servicio, denominado
evmaudit-web, se fundamenta en tres pilares de ingeniería:

\begin{itemize}
\tightlist
\item
  \textbf{Persistencia de Datos mediante Volúmenes}: Con el fin de dotar
  al sistema de un estado persistente (pese a la naturaleza efímera de
  los contenedores), se realiza un mapeo directo de directorios del host
  hacia el contenedor:
\item
  \textbf{\texttt{./jsons/\_uploads:/app/jsons/\_uploads}}: Almacena de
  forma persistente los contratos Solidity cargados por los usuarios,
  los wrappers intermedios generados para Echidna y los informes de
  auditoría finales en formato JSON y Markdown.
\item
  \textbf{\texttt{./contracts:/app/contracts}}: Habilita un volumen
  opcional para la auditoría directa de Smart Contracts locales sin
  necesidad de interactuar con la interfaz web.
\item
  \textbf{Aislamiento de Red y Mapeo de Puertos}: Se expone el puerto
  8080 del contenedor hacia el puerto 8080 del sistema anfitrión. Esto
  permite redirigir el tráfico HTTP de la interfaz construida en
  HTML5/Vanilla JS hacia el backend desarrollado en FastAPI de forma
  transparente.
\item
  \textbf{Tolerancia a Fallos}: Se implementa la política de reinicio
  restart: unless-stopped. Esta directiva asegura la alta disponibilidad
  del servicio ante excepciones imprevistas en el motor de ejecución
  simbólica (Mythril) o caídas del propio demonio de Docker,
  garantizando que el servicio web vuelva a levantarse de manera
  automática salvo interrupción explícita del administrador.
\end{itemize}

\subsection{ANEXO D.5. Flujo de Despliegue y Ciclo de Vida del
Contenedor}\label{anexo-d.5.-flujo-de-despliegue-y-ciclo-de-vida-del-contenedor}

Para la puesta en marcha de la infraestructura local en entornos de
desarrollo o evaluación, el ciclo de vida del contenedor se administra
mediante el estándar de comandos de Docker Compose: 1. \textbf{Fase de
Construcción (Build):} Compilación de la imagen e instalación del
entorno virtual determinista:

\begin{Shaded}
\begin{Highlighting}[]
\ExtensionTok{docker}\NormalTok{ compose build}
\end{Highlighting}
\end{Shaded}

\begin{enumerate}
\def\labelenumi{\arabic{enumi}.}
\setcounter{enumi}{1}
\tightlist
\item
  \textbf{Fase de Inicialización (Up):} Despliegue e instanciación del
  servicio en segundo plano (detached mode):
\end{enumerate}

\begin{Shaded}
\begin{Highlighting}[]
\ExtensionTok{docker}\NormalTok{ compose up }\AttributeTok{{-}d}
\end{Highlighting}
\end{Shaded}

\begin{enumerate}
\def\labelenumi{\arabic{enumi}.}
\setcounter{enumi}{2}
\tightlist
\item
  \textbf{Fase de Auditoría de Ejecución (Logs):} Inspección de la
  salida estándar del contenedor para la monitorización de los análisis
  en curso:
\end{enumerate}

\begin{Shaded}
\begin{Highlighting}[]
\ExtensionTok{docker}\NormalTok{ compose logs }\AttributeTok{{-}f}\NormalTok{ evmaudit{-}web}
\end{Highlighting}
\end{Shaded}

\begin{enumerate}
\def\labelenumi{\arabic{enumi}.}
\setcounter{enumi}{3}
\tightlist
\item
  \textbf{Fase de Parada (Down):} Interrupción y eliminación de los
  contenedores activos salvaguardando la integridad de los datos de las
  auditorías gracias a los volúmenes enlazados: docker compose down
\end{enumerate}
